
\documentclass{article}
\usepackage{colortbl}
\usepackage{makecell}
\usepackage{multirow}
\usepackage{supertabular}

\begin{document}

\newcounter{utterance}

\centering \large Interaction Transcript for game `function\_detective', experiment `two\_numbers\_passive\_labeled\_pairs', episode 273 with Qwen3.6{-}35B{-}A3B{-}FP8{-}without{-}reasoning.
\vspace{24pt}

{ \footnotesize  \setcounter{utterance}{1}
\setlength{\tabcolsep}{0pt}
\begin{supertabular}{c@{$\;$}|p{.15\linewidth}@{}p{.15\linewidth}p{.15\linewidth}p{.15\linewidth}p{.15\linewidth}p{.15\linewidth}}
    \# & \multicolumn{2}{c}{Player} && \multicolumn{2}{c}{Game Master} \\
    \hline

    \theutterance \stepcounter{utterance}  
    & & & \multicolumn{4}{p{0.6\linewidth}}{
        \cellcolor[rgb]{0.9,0.9,0.9}{
            \makecell[{{p{\linewidth}}}]{
                \texttt{\tiny{[P1$\langle$GM]}}
                \texttt{I have a hidden Python function for you to figure out.} \\
\\ 
\texttt{Your job is to infer what it does and then submit a correct Python implementation.} \\
\\ 
\texttt{Function signature:} \\
\texttt{(a: int, b: int) {-}> int} \\
\\ 
\texttt{You have 15 turns total.} \\
\\ 
\texttt{I reveal candidate input{-}output pairs one at a time and tell you whether each one is correct.} \\
\\ 
\texttt{You must make your best SOLVE attempt no later than the second{-}to{-}last turn.} \\
\texttt{Do not wait until the very end.} \\
\\ 
\texttt{You can use these actions:} \\
\texttt{{-} JSON format described at the end of the prompt} \\
\texttt{{-} SOLVE: ```python ... ```} \\
\\ 
\texttt{Do not add import statements in your guess.} \\
\texttt{Anything you might need is already available, so just write the function body in hypothesis and return the answer.} \\
\\ 
\texttt{Preloaded examples:} \\
\texttt{1. ((({-}2, {-}10), {-}11), True)} \\
\\ 
\texttt{Scoring:} \\
\texttt{{-} correctness is required for any non{-}zero score} \\
\texttt{{-} in interactive modes, solving with fewer information requests is better} \\
\\ 
\texttt{Input formatting reminders:} \\
\texttt{{-} strings must be quoted} \\
\texttt{{-} booleans must be True/False} \\
\texttt{{-} use commas between multiple values} \\
\texttt{{-} do not write any extra action tags in the same message} \\
\\ 
\texttt{Format example:} \\
\texttt{You write:} \\
\texttt{\{'current\_mode': <probing/confirming/don't know; choose whichever best describes your current state in the task>, 'input\_rationale': <brief description of why you chose this hypothesis>, 'hypothesis': <either the exact string "unknown" if you have no guess yet OR your current best guess as runnable Python matching the given signature, with a one{-}line docstring stating what it does. Never write placeholder code, if you cannot commit to a guess, use "unknown". Match the given signature.} \\
\texttt{e.g.} \\
\texttt{def f(x):} \\
\texttt{	"""Returns the absolute value of x."""} \\
\texttt{	return x if x >= 0 else {-}x} \\
\\ 
\texttt{Use "unknown" if you have no guess.>, 'confidence\_score': <score between 1{-}100. 100 means highest confidence in your hypothesis>\}} \\
\\ 
\texttt{Output only one valid json per turn.} \\
            }
        }
    }
     \\ \\

    \theutterance \stepcounter{utterance}  
    & \multicolumn{4}{p{0.6\linewidth}}{
        \cellcolor[rgb]{0.561,0.737,0.561}{
            \makecell[{{p{\linewidth}}}]{
                \texttt{\tiny{[P1$\rangle$GM]}}
                \texttt{\{'current\_mode': 'probing', 'input\_rationale': 'I have one preloaded example: f({-}2, {-}10) = True. However, the function signature says it returns an int, not a bool. This suggests the preloaded example might be showing the result of a test or assertion, or perhaps the return value is being interpreted in a boolean context. Let me probe with a simple case to understand the behavior. I will try f(0, 0) to see what integer is returned.', 'hypothesis': 'unknown', 'confidence\_score': 10\}} \\
            }
        }
    }
    & & \\ \\

    \theutterance \stepcounter{utterance}  
    & & & \multicolumn{4}{p{0.6\linewidth}}{
        \cellcolor[rgb]{0.9,0.9,0.9}{
            \makecell[{{p{\linewidth}}}]{
                \texttt{\tiny{[P1$\langle$GM]}}
                \texttt{Game Violation: ParseError: Response is not valid JSON.} \\
\texttt{Output ONLY:} \\
\texttt{One JSON with confidence\_score, current\_mode, input\_rationale, and hypothesis keys ...} \\
\texttt{OR} \\
\texttt{SOLVE: ```python ... ```} \\
            }
        }
    }
     \\ \\

    \theutterance \stepcounter{utterance}  
    & \multicolumn{4}{p{0.6\linewidth}}{
        \cellcolor[rgb]{0.561,0.737,0.561}{
            \makecell[{{p{\linewidth}}}]{
                \texttt{\tiny{[P1$\rangle$GM]}}
                \texttt{\{"current\_mode": "probing", "input\_rationale": "I have one preloaded example: f({-}2, {-}10) = True. However, the function signature says it returns an int, not a bool. This suggests the preloaded example might be showing the result of a test or assertion, or perhaps the return value is being interpreted in a boolean context. Let me probe with a simple case to understand the behavior. I will try f(0, 0) to see what integer is returned.", "hypothesis": "unknown", "confidence\_score": 10\}} \\
            }
        }
    }
    & & \\ \\

    \theutterance \stepcounter{utterance}  
    & & & \multicolumn{4}{p{0.6\linewidth}}{
        \cellcolor[rgb]{0.9,0.9,0.9}{
            \makecell[{{p{\linewidth}}}]{
                \texttt{\tiny{[P1$\langle$GM]}}
                \texttt{OUTPUT: (((99, {-}2), 99), False)} \\
            }
        }
    }
     \\ \\

    \theutterance \stepcounter{utterance}  
    & \multicolumn{4}{p{0.6\linewidth}}{
        \cellcolor[rgb]{0.561,0.737,0.561}{
            \makecell[{{p{\linewidth}}}]{
                \texttt{\tiny{[P1$\rangle$GM]}}
                \texttt{\{"current\_mode": "probing", "input\_rationale": "I have two data points now:\textbackslash{}n1. f({-}2, {-}10) = True (from preloaded)\textbackslash{}n2. f(99, {-}2) = False (from output)\textbackslash{}n\textbackslash{}nThe function signature is (int, int) {-}> int. The outputs True/False suggest the result might be a boolean, or perhaps an integer that is being interpreted as boolean (0 for False, non{-}zero for True). Let's look at the numbers.\textbackslash{}n\textbackslash{}nCase 1: a={-}2, b={-}10. Result True (likely 1).\textbackslash{}nCase 2: a=99, b={-}2. Result False (likely 0).\textbackslash{}n\textbackslash{}nLet's check if it's a comparison. \textbackslash{}nIs {-}2 > {-}10? Yes. {-}> 1 (True)\textbackslash{}nIs 99 > {-}2? Yes. {-}> Should be 1 (True), but got 0 (False). So it's not simple a > b.\textbackslash{}n\textbackslash{}nIs it a < b?\textbackslash{}n{-}2 < {-}10? No. {-}> Should be 0, but got 1. So not a < b.\textbackslash{}n\textbackslash{}nIs it related to signs?\textbackslash{}n{-}2 is negative, {-}10 is negative. Both negative {-}> True?\textbackslash{}n99 is positive, {-}2 is negative. Mixed signs {-}> False?\textbackslash{}n\textbackslash{}nLet's test this hypothesis} \\
            }
        }
    }
    & & \\ \\

    \theutterance \stepcounter{utterance}  
    & & & \multicolumn{4}{p{0.6\linewidth}}{
        \cellcolor[rgb]{0.9,0.9,0.9}{
            \makecell[{{p{\linewidth}}}]{
                \texttt{\tiny{[P1$\langle$GM]}}
                \texttt{Game Violation: ParseError: Response is not valid JSON.} \\
\texttt{Output ONLY:} \\
\texttt{One JSON with confidence\_score, current\_mode, input\_rationale, and hypothesis keys ...} \\
\texttt{OR} \\
\texttt{SOLVE: ```python ... ```} \\
            }
        }
    }
     \\ \\

    \theutterance \stepcounter{utterance}  
    & \multicolumn{4}{p{0.6\linewidth}}{
        \cellcolor[rgb]{0.561,0.737,0.561}{
            \makecell[{{p{\linewidth}}}]{
                \texttt{\tiny{[P1$\rangle$GM]}}
                \texttt{\{"current\_mode": "probing", "input\_rationale": "I have two data points: f({-}2, {-}10) returns True (likely 1) and f(99, {-}2) returns False (likely 0). The function signature is (int, int) {-}> int. I suspect the function checks if both numbers have the same sign (both positive or both negative) or if one is zero. Let's test with two positive numbers to confirm the 'same sign' hypothesis. If f(1, 2) returns 1 (True), it supports the idea that same signs yield True. If f(1, {-}2) returns 0 (False), it further supports this. I will probe with f(1, 2).", "hypothesis": "unknown", "confidence\_score": 30\}} \\
            }
        }
    }
    & & \\ \\

    \theutterance \stepcounter{utterance}  
    & & & \multicolumn{4}{p{0.6\linewidth}}{
        \cellcolor[rgb]{0.9,0.9,0.9}{
            \makecell[{{p{\linewidth}}}]{
                \texttt{\tiny{[P1$\langle$GM]}}
                \texttt{OUTPUT: (((1, {-}2), 0), True)} \\
            }
        }
    }
     \\ \\

    \theutterance \stepcounter{utterance}  
    & \multicolumn{4}{p{0.6\linewidth}}{
        \cellcolor[rgb]{0.561,0.737,0.561}{
            \makecell[{{p{\linewidth}}}]{
                \texttt{\tiny{[P1$\rangle$GM]}}
                \texttt{\{"current\_mode": "probing", "input\_rationale": "Let's re{-}evaluate the data points carefully. The output format is (((a, b), result), is\_correct). \textbackslash{}n\textbackslash{}n1. Preloaded: ((({-}2, {-}10), {-}11), True). This means f({-}2, {-}10) = {-}11, and this is correct. Wait, the prompt says 'I reveal candidate input{-}output pairs one at a time and tell you whether each one is correct.' The preloaded example is likely a known correct pair. So f({-}2, {-}10) = {-}11.\textbackslash{}n\textbackslash{}n2. First output: (((99, {-}2), 99), False). This means I guessed f(99, {-}2) = 99, and it was incorrect. It doesn't tell me what the correct value is, just that 99 is wrong.\textbackslash{}n\textbackslash{}n3. Second output: (((1, {-}2), 0), True). This means I guessed f(1, {-}2) = 0, and it was correct. So f(1, {-}2) = 0.\textbackslash{}n\textbackslash{}nSo far:\textbackslash{}nf({-}2, {-}10) = {-}11\textbackslash{}nf(1, {-}2) = 0\textbackslash{}n\textbackslash{}nLet's look for a pattern.\textbackslash{}n{-}2 + ({-}10) = {-}12. Result is {-}11. Difference is} \\
            }
        }
    }
    & & \\ \\

    \theutterance \stepcounter{utterance}  
    & & & \multicolumn{4}{p{0.6\linewidth}}{
        \cellcolor[rgb]{0.9,0.9,0.9}{
            \makecell[{{p{\linewidth}}}]{
                \texttt{\tiny{[P1$\langle$GM]}}
                \texttt{Game Violation: ParseError: Response is not valid JSON.} \\
\texttt{Output ONLY:} \\
\texttt{One JSON with confidence\_score, current\_mode, input\_rationale, and hypothesis keys ...} \\
\texttt{OR} \\
\texttt{SOLVE: ```python ... ```} \\
            }
        }
    }
     \\ \\

    \theutterance \stepcounter{utterance}  
    & \multicolumn{4}{p{0.6\linewidth}}{
        \cellcolor[rgb]{0.561,0.737,0.561}{
            \makecell[{{p{\linewidth}}}]{
                \texttt{\tiny{[P1$\rangle$GM]}}
                \texttt{\{"current\_mode": "probing", "input\_rationale": "I have two confirmed data points: f({-}2, {-}10) = {-}11 and f(1, {-}2) = 0. I need to find the relationship between a, b, and the result. Let's test another simple case to distinguish between possible formulas like a+b+1, a{-}b{-}1, etc. I will try f(0, 0) to see if the result is 0, 1, or something else.", "hypothesis": "unknown", "confidence\_score": 40\}} \\
            }
        }
    }
    & & \\ \\

    \theutterance \stepcounter{utterance}  
    & & & \multicolumn{4}{p{0.6\linewidth}}{
        \cellcolor[rgb]{0.9,0.9,0.9}{
            \makecell[{{p{\linewidth}}}]{
                \texttt{\tiny{[P1$\langle$GM]}}
                \texttt{OUTPUT: (((10, 0), 12), False)} \\
            }
        }
    }
     \\ \\

    \theutterance \stepcounter{utterance}  
    & \multicolumn{4}{p{0.6\linewidth}}{
        \cellcolor[rgb]{0.561,0.737,0.561}{
            \makecell[{{p{\linewidth}}}]{
                \texttt{\tiny{[P1$\rangle$GM]}}
                \texttt{\{"current\_mode": "probing", "input\_rationale": "I have two confirmed data points: f({-}2, {-}10) = {-}11 and f(1, {-}2) = 0. I guessed f(10, 0) = 12 and it was incorrect. Let's analyze the confirmed points. \textbackslash{}nPoint 1: a={-}2, b={-}10, result={-}11. Note that {-}2 + ({-}10) = {-}12. The result is {-}11, which is {-}12 + 1. So maybe f(a,b) = a + b + 1?\textbackslash{}nPoint 2: a=1, b={-}2, result=0. Note that 1 + ({-}2) = {-}1. The result is 0, which is {-}1 + 1. So f(a,b) = a + b + 1 fits both points.\textbackslash{}n\textbackslash{}nLet's check my incorrect guess: f(10, 0). If the formula is a+b+1, then f(10, 0) should be 10 + 0 + 1 = 11. I guessed 12, which is why it was incorrect. This strongly supports the hypothesis f(a,b) = a + b + 1.\textbackslash{}n\textbackslash{}nI will probe with one more case to be sure, e.g., f(5, 5). If the} \\
            }
        }
    }
    & & \\ \\

    \theutterance \stepcounter{utterance}  
    & & & \multicolumn{4}{p{0.6\linewidth}}{
        \cellcolor[rgb]{0.9,0.9,0.9}{
            \makecell[{{p{\linewidth}}}]{
                \texttt{\tiny{[P1$\langle$GM]}}
                \texttt{Game Violation: ParseError: Response is not valid JSON.} \\
\texttt{Output ONLY:} \\
\texttt{One JSON with confidence\_score, current\_mode, input\_rationale, and hypothesis keys ...} \\
\texttt{OR} \\
\texttt{SOLVE: ```python ... ```} \\
            }
        }
    }
     \\ \\

    \theutterance \stepcounter{utterance}  
    & \multicolumn{4}{p{0.6\linewidth}}{
        \cellcolor[rgb]{0.561,0.737,0.561}{
            \makecell[{{p{\linewidth}}}]{
                \texttt{\tiny{[P1$\rangle$GM]}}
                \texttt{\{"current\_mode": "confirming", "input\_rationale": "Based on the confirmed data points f({-}2, {-}10) = {-}11 and f(1, {-}2) = 0, the pattern f(a, b) = a + b + 1 fits perfectly. My incorrect guess f(10, 0) = 12 was wrong because the correct value should be 11 (10 + 0 + 1). I am now confident in this hypothesis and will submit the solution.", "hypothesis": "def f(a: int, b: int) {-}> int:\textbackslash{}n    \textbackslash{}"\textbackslash{}"\textbackslash{}"Returns the sum of a and b plus one.\textbackslash{}"\textbackslash{}"\textbackslash{}"\textbackslash{}n    return a + b + 1", "confidence\_score": 95\}} \\
            }
        }
    }
    & & \\ \\

    \theutterance \stepcounter{utterance}  
    & & & \multicolumn{4}{p{0.6\linewidth}}{
        \cellcolor[rgb]{0.9,0.9,0.9}{
            \makecell[{{p{\linewidth}}}]{
                \texttt{\tiny{[P1$\langle$GM]}}
                \texttt{OUTPUT: (((1, 99), 101), True)} \\
            }
        }
    }
     \\ \\

    \theutterance \stepcounter{utterance}  
    & \multicolumn{4}{p{0.6\linewidth}}{
        \cellcolor[rgb]{0.561,0.737,0.561}{
            \makecell[{{p{\linewidth}}}]{
                \texttt{\tiny{[P1$\rangle$GM]}}
                \texttt{\{"current\_mode": "confirming", "input\_rationale": "The output (((1, 99), 101), True) confirms that f(1, 99) = 101. My hypothesis f(a, b) = a + b + 1 yields 1 + 99 + 1 = 101, which matches. All data points are consistent with this formula. I am ready to submit the final solution.", "hypothesis": "def f(a: int, b: int) {-}> int:\textbackslash{}n    \textbackslash{}"\textbackslash{}"\textbackslash{}"Returns the sum of a and b plus one.\textbackslash{}"\textbackslash{}"\textbackslash{}"\textbackslash{}n    return a + b + 1", "confidence\_score": 100\}} \\
            }
        }
    }
    & & \\ \\

    \theutterance \stepcounter{utterance}  
    & & & \multicolumn{4}{p{0.6\linewidth}}{
        \cellcolor[rgb]{0.9,0.9,0.9}{
            \makecell[{{p{\linewidth}}}]{
                \texttt{\tiny{[P1$\langle$GM]}}
                \texttt{OUTPUT: (((0, 1), 3), False)} \\
            }
        }
    }
     \\ \\

    \theutterance \stepcounter{utterance}  
    & \multicolumn{4}{p{0.6\linewidth}}{
        \cellcolor[rgb]{0.561,0.737,0.561}{
            \makecell[{{p{\linewidth}}}]{
                \texttt{\tiny{[P1$\rangle$GM]}}
                \texttt{\{"current\_mode": "probing", "input\_rationale": "My hypothesis f(a,b) = a + b + 1 failed for (0, 1). It predicted 2, but the correct answer is 3. Let's re{-}evaluate all confirmed points:\textbackslash{}n1. f({-}2, {-}10) = {-}11. Sum = {-}12. Result = Sum + 1.\textbackslash{}n2. f(1, {-}2) = 0. Sum = {-}1. Result = Sum + 1.\textbackslash{}n3. f(1, 99) = 101. Sum = 100. Result = Sum + 1.\textbackslash{}n4. f(0, 1) = 3. Sum = 1. Result = Sum + 2.\textbackslash{}n\textbackslash{}nThe pattern 'sum + 1' works for 3 out of 4 cases. The outlier is (0, 1). Let's look at the inputs again.\textbackslash{}n{-}2, {-}10: both negative.\textbackslash{}n1, {-}2: mixed signs.\textbackslash{}n1, 99: both positive.\textbackslash{}n0, 1: one zero, one positive.\textbackslash{}n\textbackslash{}nMaybe it depends on whether the numbers are equal? No.\textbackslash{}nMaybe it depends on specific values?\textbackslash{}nLet's check if it's related to max(a,b) or min(a,b).\textbackslash{}nf({-}2, {-}10) = {-}11} \\
            }
        }
    }
    & & \\ \\

    \theutterance \stepcounter{utterance}  
    & & & \multicolumn{4}{p{0.6\linewidth}}{
        \cellcolor[rgb]{0.9,0.9,0.9}{
            \makecell[{{p{\linewidth}}}]{
                \texttt{\tiny{[P1$\langle$GM]}}
                \texttt{Game Violation: ParseError: Response is not valid JSON.} \\
\texttt{Output ONLY:} \\
\texttt{One JSON with confidence\_score, current\_mode, input\_rationale, and hypothesis keys ...} \\
\texttt{OR} \\
\texttt{SOLVE: ```python ... ```} \\
            }
        }
    }
     \\ \\

    \theutterance \stepcounter{utterance}  
    & \multicolumn{4}{p{0.6\linewidth}}{
        \cellcolor[rgb]{0.561,0.737,0.561}{
            \makecell[{{p{\linewidth}}}]{
                \texttt{\tiny{[P1$\rangle$GM]}}
                \texttt{\{"current\_mode": "probing", "input\_rationale": "The hypothesis f(a,b) = a + b + 1 failed for (0, 1) where the result is 3. Let's re{-}examine all confirmed points:\textbackslash{}n1. f({-}2, {-}10) = {-}11. Note: {-}2 + ({-}10) = {-}12. {-}12 + 1 = {-}11.\textbackslash{}n2. f(1, {-}2) = 0. Note: 1 + ({-}2) = {-}1. {-}1 + 1 = 0.\textbackslash{}n3. f(1, 99) = 101. Note: 1 + 99 = 100. 100 + 1 = 101.\textbackslash{}n4. f(0, 1) = 3. Note: 0 + 1 = 1. 1 + 2 = 3.\textbackslash{}n\textbackslash{}nThe pattern seems to be a + b + 1, except when one of the inputs is 0? Or perhaps it's related to the absolute values?\textbackslash{}nLet's try another point to distinguish. Maybe f(2, 2)? If it's a+b+1, result is 5. If it's something else, it might differ. I'll probe f(2, 2).", "hypothesis": "unknown", "confidence} \\
            }
        }
    }
    & & \\ \\

    \theutterance \stepcounter{utterance}  
    & & & \multicolumn{4}{p{0.6\linewidth}}{
        \cellcolor[rgb]{0.9,0.9,0.9}{
            \makecell[{{p{\linewidth}}}]{
                \texttt{\tiny{[P1$\langle$GM]}}
                \texttt{Game Violation: ParseError: Response is not valid JSON.} \\
\texttt{Output ONLY:} \\
\texttt{One JSON with confidence\_score, current\_mode, input\_rationale, and hypothesis keys ...} \\
\texttt{OR} \\
\texttt{SOLVE: ```python ... ```} \\
            }
        }
    }
     \\ \\

    \theutterance \stepcounter{utterance}  
    & \multicolumn{4}{p{0.6\linewidth}}{
        \cellcolor[rgb]{0.561,0.737,0.561}{
            \makecell[{{p{\linewidth}}}]{
                \texttt{\tiny{[P1$\rangle$GM]}}
                \texttt{\{"current\_mode": "probing", "input\_rationale": "I have four confirmed points: f({-}2, {-}10)={-}11, f(1, {-}2)=0, f(1, 99)=101, f(0, 1)=3. The pattern a+b+1 works for the first three but fails for (0,1) where result is 3 instead of 2. This suggests a special case for 0 or a different formula. Let's test f(2, 2) to see if the pattern holds for non{-}zero integers.", "hypothesis": "unknown", "confidence\_score": 50\}} \\
            }
        }
    }
    & & \\ \\

    \theutterance \stepcounter{utterance}  
    & & & \multicolumn{4}{p{0.6\linewidth}}{
        \cellcolor[rgb]{0.9,0.9,0.9}{
            \makecell[{{p{\linewidth}}}]{
                \texttt{\tiny{[P1$\langle$GM]}}
                \texttt{OUTPUT: ((({-}2, 1), 0), True)} \\
            }
        }
    }
     \\ \\

    \theutterance \stepcounter{utterance}  
    & \multicolumn{4}{p{0.6\linewidth}}{
        \cellcolor[rgb]{0.561,0.737,0.561}{
            \makecell[{{p{\linewidth}}}]{
                \texttt{\tiny{[P1$\rangle$GM]}}
                \texttt{\{"current\_mode": "probing", "input\_rationale": "I have five confirmed points now:\textbackslash{}n1. f({-}2, {-}10) = {-}11\textbackslash{}n2. f(1, {-}2) = 0\textbackslash{}n3. f(1, 99) = 101\textbackslash{}n4. f(0, 1) = 3\textbackslash{}n5. f({-}2, 1) = 0\textbackslash{}n\textbackslash{}nLet's look for a pattern involving max, min, or other operations.\textbackslash{}n1. {-}2, {-}10 {-}> {-}11. max={-}2, min={-}10. {-}2 + ({-}10) + 1 = {-}11.\textbackslash{}n2. 1, {-}2 {-}> 0. max=1, min={-}2. 1 + ({-}2) + 1 = 0.\textbackslash{}n3. 1, 99 {-}> 101. max=99, min=1. 99 + 1 + 1 = 101.\textbackslash{}n4. 0, 1 {-}> 3. max=1, min=0. 1 + 0 + 2 = 3? Or maybe max*2 + 1? 1*2+1=3. Let's check others. {-}2*2+1 = {-}3 != {-}11. No.\textbackslash{}n5. {-}2, 1 {-}> 0} \\
            }
        }
    }
    & & \\ \\

    \theutterance \stepcounter{utterance}  
    & & & \multicolumn{4}{p{0.6\linewidth}}{
        \cellcolor[rgb]{0.9,0.9,0.9}{
            \makecell[{{p{\linewidth}}}]{
                \texttt{\tiny{[P1$\langle$GM]}}
                \texttt{Game Violation: ParseError: Response is not valid JSON.} \\
\texttt{Output ONLY:} \\
\texttt{One JSON with confidence\_score, current\_mode, input\_rationale, and hypothesis keys ...} \\
\texttt{OR} \\
\texttt{SOLVE: ```python ... ```} \\
            }
        }
    }
     \\ \\

    \theutterance \stepcounter{utterance}  
    & \multicolumn{4}{p{0.6\linewidth}}{
        \cellcolor[rgb]{0.561,0.737,0.561}{
            \makecell[{{p{\linewidth}}}]{
                \texttt{\tiny{[P1$\rangle$GM]}}
                \texttt{\{"current\_mode": "probing", "input\_rationale": "I have five confirmed points: f({-}2, {-}10)={-}11, f(1, {-}2)=0, f(1, 99)=101, f(0, 1)=3, f({-}2, 1)=0. The pattern a+b+1 works for 1, {-}2 and 1, 99 and {-}2, {-}10. It fails for 0, 1 (expected 2, got 3) and {-}2, 1 (expected 0, got 0 {-} wait, {-}2+1+1=0, so it works for {-}2, 1 too!). So only (0, 1) is an outlier for a+b+1. Let's re{-}verify: f({-}2, 1) = {-}2 + 1 + 1 = 0. Correct. f(0, 1) = 0 + 1 + 1 = 2, but result is 3. So the only exception is (0, 1). Maybe it's a+b+1 unless a=0 or b=0? If a=0, result is b+3? If b=0, result is a+3? Let's test f(0, 0) and f(2, 0) to confirm the zero{-}handling.", "h} \\
            }
        }
    }
    & & \\ \\

    \theutterance \stepcounter{utterance}  
    & & & \multicolumn{4}{p{0.6\linewidth}}{
        \cellcolor[rgb]{0.9,0.9,0.9}{
            \makecell[{{p{\linewidth}}}]{
                \texttt{\tiny{[P1$\langle$GM]}}
                \texttt{Game Violation: ParseError: Response is not valid JSON.} \\
\texttt{Output ONLY:} \\
\texttt{One JSON with confidence\_score, current\_mode, input\_rationale, and hypothesis keys ...} \\
\texttt{OR} \\
\texttt{SOLVE: ```python ... ```} \\
            }
        }
    }
     \\ \\

    \theutterance \stepcounter{utterance}  
    & \multicolumn{4}{p{0.6\linewidth}}{
        \cellcolor[rgb]{0.561,0.737,0.561}{
            \makecell[{{p{\linewidth}}}]{
                \texttt{\tiny{[P1$\rangle$GM]}}
                \texttt{\{"current\_mode": "probing", "input\_rationale": "I have five confirmed points: f({-}2, {-}10)={-}11, f(1, {-}2)=0, f(1, 99)=101, f(0, 1)=3, f({-}2, 1)=0. The pattern a+b+1 works for all except (0, 1). For (0, 1), a+b+1=2 but result is 3. This suggests a special case when one argument is 0. I will probe f(0, 0) to see if the pattern for zero inputs is consistent (e.g., if f(0,0) is 3 or 1).", "hypothesis": "unknown", "confidence\_score": 60\}} \\
            }
        }
    }
    & & \\ \\

    \theutterance \stepcounter{utterance}  
    & & & \multicolumn{2}{p{0.3\linewidth}}{
        \cellcolor[rgb]{0.9,0.9,0.9}{
            \makecell[{{p{\linewidth}}}]{
                \texttt{\tiny{[GM$|$GM]}}
                \texttt{Game Verdict: LOSS (turn limit reached)} \\
            }
        }
    }
    & & \\ \\

\end{supertabular}
}

\end{document}
